% ================================================================
% CONFIGURATION GLOBALE : éditer principalement les lignes ci-dessous
% Compatible Overleaf gratuit / compilateur pdfLaTeX / sans options payantes
% ================================================================
\usepackage[T1]{fontenc}
\usepackage[utf8]{inputenc}
\usepackage[french]{babel}
\usepackage{lmodern}
\usepackage{microtype}
\usepackage[a4paper,left=27mm,right=25mm,top=27mm,bottom=26mm,headheight=16pt]{geometry}
\usepackage{graphicx}
\usepackage{tikz}
\usetikzlibrary{calc,positioning,arrows.meta}
\usepackage[most]{tcolorbox}
\usepackage{amsmath,amssymb}
\usepackage{booktabs,tabularx,array,longtable}
\usepackage{enumitem}
\usepackage{titlesec}
\usepackage{fancyhdr}
\usepackage{lastpage}
\usepackage{caption}
\usepackage{float}
\usepackage{hyperref}
\usepackage{url}
\usepackage{xcolor}

% --- Identité du document : MODIFIER CES CHAMPS ---
\newcommand{\Ecole}{École Centrale Casablanca}
\newcommand{\Cours}{Machine Learning et Intelligence Artificielle}
\newcommand{\TitreRapport}{Autoencodeurs pour le clustering\\et la détection d'anomalies}
\newcommand{\SousTitreRapport}{Application à la surveillance vibratoire des roulements}
\newcommand{\MembreA}{Prénom NOM du membre A}
\newcommand{\MembreB}{Prénom NOM du membre B}
\newcommand{\MembreC}{Prénom NOM du membre C}
\newcommand{\MembreD}{Prénom NOM du membre D}
\newcommand{\Encadrant}{À compléter}
\newcommand{\AnneeProjet}{2026}
\newcommand{\PeriodeProjet}{22 octobre -- 30 novembre 2026}
\newcommand{\VersionRapport}{0.1}
\newcommand{\DateMiseAJour}{27 octobre 2026}

% --- Logos et photo FACULTATIFS ---
% Ajouter ces fichiers dans le dossier images/, sans toucher au code.
\newcommand{\CheminLogoECC}{images/logo_ecc.png}
\newcommand{\CheminLogoPartenaire}{images/logo_partenaire.png}
\newcommand{\CheminPhotoBanc}{images/banc_essai.jpg}

% --- Palette sobre, académique, industrielle ---
\definecolor{bleunuit}{HTML}{14253D}
\definecolor{bleufonce}{HTML}{223C59}
\definecolor{turquoise}{HTML}{0E8794}
\definecolor{orclair}{HTML}{DBAD54}
\definecolor{grisclair}{HTML}{F3F6F8}
\definecolor{grisfonce}{HTML}{596A77}
\definecolor{bleupale}{HTML}{EAF3F5}
\definecolor{orangealerte}{HTML}{A36B21}

\hypersetup{pdftitle={Autoencodeurs pour le clustering et la détection d'anomalies},
            pdfauthor={Équipe de projet, École Centrale Casablanca},
            colorlinks=true,linkcolor=bleufonce,citecolor=turquoise,urlcolor=turquoise}
\setlength{\parindent}{0pt}
\setlength{\parskip}{0.55em}
\renewcommand{\arraystretch}{1.18}
\setlist[itemize]{leftmargin=*,topsep=3pt,itemsep=2pt}
\setlist[enumerate]{leftmargin=*,topsep=3pt,itemsep=2pt}
\captionsetup{font=small,labelfont={bf,color=bleunuit}}

% --- Titres et navigation ---
\titleformat{\chapter}[display]
  {\sffamily\bfseries\color{bleunuit}}
  {\textcolor{turquoise}{\small\MakeUppercase{Chapitre}\ \thechapter}}
  {3pt}{\Huge}
\titlespacing*{\chapter}{0pt}{-6pt}{20pt}
\titleformat{\section}{\Large\sffamily\bfseries\color{bleufonce}}{\thesection}{0.8em}{}
\titleformat{\subsection}{\normalsize\sffamily\bfseries\color{turquoise}}{\thesubsection}{0.8em}{}

\pagestyle{fancy}
\fancyhf{}
\fancyhead[L]{\small\sffamily\color{grisfonce}Projet Autoencodeurs $\cdot$ ML \& IA}
\fancyhead[R]{\small\sffamily\color{turquoise}Document évolutif $\vert$ v\VersionRapport}
\fancyfoot[L]{\footnotesize\color{grisfonce}\Ecole}
\fancyfoot[R]{\footnotesize\color{grisfonce}\thepage}
\renewcommand{\headrulewidth}{0.3pt}
\renewcommand{\footrulewidth}{0pt}
\fancypagestyle{plain}{\fancyhf{}\fancyfoot[L]{\footnotesize\color{grisfonce}\Ecole}\fancyfoot[R]{\footnotesize\color{grisfonce}\thepage}\renewcommand{\headrulewidth}{0pt}}

% --- Blocs d'information réutilisables ---
\newtcolorbox{encadreinfo}[1][]{enhanced,breakable,colback=bleupale,
  colframe=turquoise!45,boxrule=0.35pt,arc=1.5mm,
  borderline west={2.5pt}{0pt}{turquoise},left=4mm,right=4mm,top=3mm,bottom=3mm,#1}
\newtcolorbox{encadreattente}[1][]{enhanced,breakable,colback=grisclair,
  colframe=grisfonce!20,boxrule=0.35pt,arc=1mm,
  borderline west={2pt}{0pt}{orclair},left=4mm,right=4mm,top=3mm,bottom=3mm,#1}
\newcommand{\acompleter}[1]{\begin{encadreattente}\textbf{À compléter au cours du projet.}\\ #1\end{encadreattente}}
